
\section{辛四维流形}\label{sec:9-4}

本节证明辛四维流形上辛结构的若干经典唯一性结果。相关问题还有很多，其中一些在
注~\ref{rem:9-4-3} 中讨论，另一些留到第~\ref{sec:9-7} 节。

\personlater{geluomofu}{格罗莫夫}{Gromov}在 \srcite{160} 中证明：若非紧辛四维流形
$(M,\omega)$ 在无穷远处辛同胚于 $(\mathbb R^4,\omega_0)$，并且 $\omega$ 在
$\pi_2(M)$ 上消失，则 $M$ 必辛同胚于 $\mathbb R^4$。他还证明了
$S^2\times S^2$ 上\BookRecordIndex{T0363}辛形式的一个唯一性结果：设\BookRecordIndex{T0363}辛形式 $\omega$ 在两个球面
\[
  C_A:=S^2\times\{\mathrm{pt}\},
  \qquad
  C_B:=\{\mathrm{pt}\}\times S^2
\]
上非退化；若这两个球面有相同面积，则 $\omega$ 微分同胚于标准乘积形式
$\pi_1^*\sigma+\pi_2^*\sigma$。他由此\BookRecordIndex{T0326}推出下面定理~\ref{thm:9-4-1}(iii) 的唯一性陈述。
McDuff \srcite{258} 在他的思想基础上证明了有理曲面和直纹曲面的如下辛分类定理。回顾：
所谓\emph{辛直纹曲面}，是指具有如下局部平凡\BookRecordIndex{T0279}纤维化结构的辛四维流形 $(M,\omega)$：
\[
  \pi:M\longrightarrow\Sigma,
\]
其中 $\Sigma$ 是黎曼曲面，各纤维都是二维球面，且 $\omega$ 在其上非退化。

\begin{theorem}[McDuff]\label{thm:9-4-1}
设 $(M,\omega)$ 是闭连通辛四维流形，它含有一个自交数非负的辛嵌入二维球面，却不含
自交数为 $-1$ 的辛嵌入二维球面。则：
\begin{enumerate}[label=(\roman*),leftmargin=3em]
\item 存在自交数为零或一的辛嵌入二维球面 $C$；
\item 若存在自交数 $C\mathbin{\boldsymbol\cdot}C=0$ 的辛嵌入二维球面 $C$，则
$M$ 辛同胚于一个直纹曲面；
\item 若存在自交数 $C\mathbin{\boldsymbol\cdot}C=1$ 的辛嵌入二维球面 $C$，则流形对
$(M,C)$ 与 $(\mathbb{CP}^2,\mathbb{CP}^1)$ 辛同胚。
\end{enumerate}
\end{theorem}

格罗莫夫的基本思想很简单。他证明，每个球面 $C_A,C_B$ 都属于一族球面，而且 $M$ 的
每个点都通过该族中的一个球面。此外，每个 $A$-球面都和每个 $B$-球面恰好横截相交一次。
由此可知这些球面都是嵌入的，因而给出 $M$ 上的新坐标。最后一步是证明，在这些新坐标
下，$\omega$ 是乘积形式。定理~\ref{thm:9-4-7} 将完成这个证明。第一步和最后一步相当
直接，中间一步却要用到定理~\ref{thm:2-6-3} 的\BookRecordIndex{T0186}交点正性，证明起来困难得多。（遗憾的是，
习题~\ref{ex:2-6-1} 描述的简单情形并不足够，因为还必须排除在两条曲线上都为奇异点的
交点。）定理~\ref{thm:9-4-1} 还需要更多准备。除了 \BookRecordIndex{T0003}$J$-全纯曲线的一般性质外，证明还
依赖 McDuff 的\BookRecordIndex{T0055}伴随公式（定理~\ref{thm:2-6-4}）以及 McDuff \srcite{257} 中关于辛吹起
和吹下的结果。下面的证明大体遵循 Lalonde--McDuff \srcite{226} 的叙述。不过，其中
一些更具几何性的论证（例如 \srcite{257} 的引理 3.5）在此由分析论证替代，后者使用实
线性\BookRecordIndex{T0210}柯西--黎曼算子的性质。

把定理~\ref{thm:9-4-1} 与格罗莫夫 \srcite{160} 的定理 0.3.C 结合起来，得到如下结果。
若 $V\subset\mathbb R^m$ 满足：$x\in V$ 时，对每个 $t\in[0,1]$ 都有 $tx\in V$，就称
$V$ 为\emph{星形的}。

\begin{theorem}[Gromov--McDuff]\label{thm:9-4-2}
设 $(M,\omega)$ 是连通辛四维流形，$K\subset M$ 是紧子集，并且满足：
\begin{enumerate}[label=(\roman*),leftmargin=3em]
\item $M$ 中不存在自交数
$S\mathbin{\boldsymbol\cdot}S=-1$ 的辛嵌入二维球面 $S$；
\item 存在辛同胚
$\psi:\mathbb R^4\setminus V\to M\setminus K$，其中
$V\subset\mathbb R^4$ 是星形紧集。
\end{enumerate}
则 $(M,\omega)$ 辛同胚于 $(\mathbb R^4,\omega_0)$。此外，对 $K$ 的每个开邻域
$U\subset M$，可以选取该辛同胚，使它在 $M\setminus U$ 上等于 $\psi^{-1}$。
\end{theorem}

在证明定理~\ref{thm:9-4-1} 和定理~\ref{thm:9-4-2} 之前，先讨论若干相关进展，并把
这些结果置于其背景中。

\begin{remark}\label{rem:9-4-3}
\begin{enumerate}[label=(\roman*),leftmargin=3em]
\item 设 $M\to\Sigma$ 是直纹曲面。McDuff \srcite{258} 在 $\Sigma=S^2$ 的情形、
Lalonde--McDuff \srcite{225} 在一般情形证明：若 $M$ 上两个上同调的\BookRecordIndex{T0363}辛形式都容许纤维
类中的辛嵌入球面，则它们微分同胚。此外，可选取微分同胚使这两个纤维类中的辛球面彼此
对应。例如，若丛 $M\to S^2$ 平凡，则 $\omega$ 微分同胚于乘积形式
$\pi_1^*\sigma_1+\pi_2^*\sigma_2$。这项工作可以理解为定理~\ref{thm:9-4-1}(ii) 的
加强。它用到一种称为\emph{\BookRecordIndex{T0359}辛膨胀}的技术，以控制\BookRecordIndex{T0363}辛形式上同调类的变化。（见
命题~\ref{prop:9-7-2} 的证明；这个方法的更多应用见 McDuff \srcite{266} 和 Biran
\srcite{38,39}。）在乘积 $M=S^2\times S^2$ 的情形，只要假设另有一个辛嵌入球面，就
可以不使用膨胀。这里需要的论证见定理~\ref{thm:9-4-7}；除了当两个球面的大小不同时，
$A,B$-曲线可能退化所造成的技术细节外，这些论证本质上出自格罗莫夫。它们也是
定理~\ref{thm:9-4-1}(iii) 证明的基础。

\item Taubes 在 \srcite{392,393} 中证明：对 $\mathbb{CP}^2$ 上每个\BookRecordIndex{T0363}辛形式 $\omega$ 和
每个\BookRecordIndex{T0088}殆复结构 $J\in\mathcal J(\mathbb{CP}^2,\omega)$，都有一条 $J$-全纯球面表示
$H_2(\mathbb{CP}^2;\mathbb Z)$ 中的 $\omega$-正生成元。把这个结果与
定理~\ref{thm:9-4-1}(iii) 结合，可知 $\mathbb{CP}^2$ 上每个体积为 $\pi^2/2$ 的\BookRecordIndex{T0363}辛形式
都辛同胚于 Fubini--Study 形式 $\omega_{\mathrm{FS}}$。注意，反向形式
$-\omega_{\mathrm{FS}}$ 经由保定向微分同胚
\[
  f:\mathbb{CP}^2\to\mathbb{CP}^2,
  \qquad
  f([z_0:z_1:z_2]):=[\bar z_0:\bar z_1:\bar z_2]
\]
辛同胚于 $\omega_{\mathrm{FS}}$。

类似地，Li--Liu \srcite{240} 利用 Taubes--Seiberg--Witten 理论证明：直纹曲面 $M$ 上
的每个\BookRecordIndex{T0363}辛形式都容许纤维类中的嵌入辛球面。结合上面讨论过的
Lalonde--McDuff \srcite{258,225,226} 的工作，可知 $M$ 上任意两个上同调的\BookRecordIndex{T0363}辛形式都
微分同胚。

\item $(M^4,\omega)$ 中自交数为 $-1$ 的辛嵌入二维球面称为\emph{例外球面}。若
$(M,\omega)$ 不含例外球面，就称它\emph{极小}。McDuff 在 \srcite{258} 中证明：若在
任意辛四维流形 $(M,\omega)$ 中吹下一族极大的两两不交例外球面，则所得\BookRecordIndex{T0358}辛流形
$(\overline M,\overline\omega)$ 是极小的。此外，她在 \srcite{260} 中利用若干曲线
\BookRecordIndex{T0246}模空间的拓扑性质证明：除非 $(\overline M,\overline\omega)$ 是有理曲面或直纹曲面，
吹下的辛同胚型与所有选择无关。如今，这个结果已经推广到光滑范畴，并且最好视作
Taubes--Seiberg--Witten 理论的推论；比较 Li--Liu \srcite{240}、Li \srcite{241}。

\item 若存在微分同胚 $f:M\to M'$ 和光滑\BookRecordIndex{T0363}辛形式族
\[
  [0,1]\longrightarrow\Omega^2(M),
  \qquad t\longmapsto\omega_t,
\]
其中各 $\omega_t$ 未必属于同一上同调类，并满足
$\omega_0=\omega$、$\omega_1=f^*\omega'$，就称两个\BookRecordIndex{T0358}辛流形
$(M,\omega)$ 与 $(M',\omega')$ \emph{形变等价}\footnote{有些作者把这个关系称为
等价，另一些作者则称之为弱形变等价。}，也称它们\emph{在形变意义下辛同胚}。现在已经
知道：对任意两个整数 $N\ge2$ 和 $k\ge1$，都存在一个 $4k$-维流形 $M$ 及其上的\BookRecordIndex{T0363}辛形式
$\omega_1,\ldots,\omega_N$，使相应\BookRecordIndex{T0358}辛流形
$(M,\omega_1),\ldots,(M,\omega_N)$ 两两不形变等价。McMullen--Taubes
\srcite{284} 给出了 $N=2,k=1$ 的首批例子，Smith \srcite{382} 则证明了上述一般结果。
他的流形由 $\mathcal J(M,\omega_i)$ 中元素的第一 Chern 类 $c_1(J)$ 的可除性区分。
与之相反，以上 (i)、(ii) 中的结果蕴含：直纹曲面上的任意两个\BookRecordIndex{T0363}辛形式在形变意义下辛同胚。

\item 设 $M$ 是闭定向 $2n$-维流形，$a\in H^2(M;\mathbb R)$ 是可由\BookRecordIndex{T0363}辛形式表示的
上同调类。用 $\Omega_a\subset\Omega^2(M)$ 表示所有代表类 $a$ 的\BookRecordIndex{T0363}辛形式所成的空间。
一个自然问题是：这个空间是否连通？若答案肯定，则 Moser 同痕说明代表类 $a$ 的任意两个
\BookRecordIndex{T0363}辛形式都微分同胚。在二维，$\Omega_a$ 总是凸的，因而连通。在六维，McDuff 给出了
$\Omega_a$ 不连通的例子（见第~\ref{sec:9-7} 节）。然而在四维，这个问题无论朝哪个方向都完全
未解：没有任何一个光滑四维流形 $M$ 和上同调类 $a\in H^2(M;\mathbb R)$，使
$\Omega_a$ 已知不连通；也没有任何一个例子，使 $\Omega_a$ 非空且已知连通。

当 $M=\mathbb{CP}^2$ 时，由以上以及第~\ref{sec:9-5} 节所述格罗莫夫、McDuff 和 Taubes 的结果，
$\Omega_a$ 连通，当且仅当 $\mathbb{CP}^2$ 上每个在同调上诱导恒等的微分同胚都同痕于
恒等映射。反过来，可以问同一上同调类 $a$ 中的两个\BookRecordIndex{T0363}辛形式能否有不同的 Chern 类；若能，
当然会\BookRecordIndex{T0326}推出 $\Omega_a$ 不连通。对 $\mathbb{CP}^2$ 的任意有限点吹起，已知同一上同调类中
的所有\BookRecordIndex{T0363}辛形式也有相同的第一 Chern 类 \srcite{240}。在 (iv) 所述 Smith 的例子中，尚不
知道那些第一 Chern 类不同的\BookRecordIndex{T0363}辛形式能否选在同一上同调类中。

\item 辛拓扑中有一个长期猜想：四维环面上的每个\BookRecordIndex{T0363}辛形式都能在同一上同调类中形变为常系数
形式。这会肯定回答 (v) 中提出的唯一性问题。写作本书时，这仍是公开问题。Donaldson
\srcite{86} 提议用一个新的几何流方程来研究这个猜想，以及 K3 曲面上的等价猜想，这一
途径颇有希望。（另见 Tosatti--Weinkove \srcite{397} 及其中的参考文献。）Taubes
\srcite{390,391} 基于 Seiberg--Witten 理论证明：四维环面或 K3 曲面上的每个\BookRecordIndex{T0363}辛形式都
有零第一 Chern 类。在 K3 曲面的情形，Donaldson 以这个论证为基础并使用上同调定向作了
改进，证明 K3 曲面上的任意两个\BookRecordIndex{T0363}辛形式（与标准定向\BookRecordIndex{T0348}相容）都属于同一个非退化二形式的同伦
类。即使这个较弱的问题，对 $T^4$ 也仍未解决。（第一 Chern 类为零的非退化二形式的同伦
类，由 $H^3(M;\mathbb Z)\oplus\mathbb Z_2$ 中的一个不变量分类。）

\item 考察辛四维流形和 凯勒 曲面之间的关系很有意思。有两个自然问题：
\begin{enumerate}[label=(\Roman*),leftmargin=3em]
\item 给定闭辛四维流形 $(M,\omega)$，是否存在与 $\omega$ \BookRecordIndex{T0348}相容的可积复结构 $J$？
\item 给定闭辛四维流形 $(M,\omega)$，$M$ 上是否存在一个与 $\omega$ 表示同一上同调类的
凯勒 形式？
\end{enumerate}
由定理~\ref{thm:9-4-1} 和 (ii) 中提到的唯一性结果，对 $\mathbb{CP}^2$ 和直纹曲面，
问题 (I) 的答案是肯定的。Kodaira 的复曲面分类定理蕴含：四维环面上的每个 凯勒 形式
都微分同胚于线性形式。比 (vi) 中问题稍弱的是，证明四维环面上的每个\BookRecordIndex{T0363}辛形式都微分同胚于
线性形式。由 Kodaira 分类，这等价于肯定回答四维环面的 (I)。

更弱的问题 (II) 涉及流形 $M$ 的\emph{\BookRecordIndex{T0367}辛锥}（即有辛代表元的上同调类所成的集合）与
\emph{凯勒 锥}（即由 $M$ 上某个复结构的 凯勒 形式表示的类所成的集合）之间的关系。
Thurston 的一个经典例子说明，\BookRecordIndex{T0358}辛流形的 凯勒 锥可以为空；这时当然对每个\BookRecordIndex{T0363}辛形式
$\omega$，问题 (I) 的答案都是否定的。如今已有许多这样的例子，其中也包括单连通例子。
不过，若两个锥都非空，已知它们不同的例子却很少。即使对
$\mathbb{CP}^2$ 的九点或更多点吹起，这个问题的答案也未知；见 Biran 关于 Nagata 猜想
的工作 \srcite{40}。然而，当 $M$ 是四维环面的单点吹起时（见
Latschev--McDuff--Schlenk \srcite{231}，推论 1.3），或当 $M$ 是二维环面上非平凡
$S^2$-丛的单点吹起时（见 Cascini--Panov \srcite{60}），这两个锥确实不同。

Tian-Jun Li 在 \srcite{243} 中提出的另一个相关问题，是理解 Kodaira 维数为零的辛
四维流形，即满足
\[
  \omega\mathbin{\boldsymbol\cdot}c_1(TM)=c_1(TM)^2=0
\]
的辛四维流形。目前已知的这类流形只有 K3 曲面、Enriques 曲面、超\BookRecordIndex{T0328}椭圆曲面和
$T^2$ 上的 $T^2$-丛。（其中包括 凯勒 锥为空的 Thurston 例子。）尽管尚不清楚这份
清单是否完备，Bauer \srcite{32} 和 Li \srcite{244} 已证明，每个这样的流形都和清单中
某个成员有相同的上同调环。关于环面丛的进一步资料另见 Smith \srcite{383}，类似问题的
讨论见 Li \srcite{245}。Streets--Tian 在 \srcite{386} 中研究的新流，即 Kähler--Ricci
流的一个辛版本，也许最终能为其中一些问题带来启发。

\item 最后要指出，本节的方法和结果是四维特有的：证明大量使用\BookRecordIndex{T0055}伴随公式，而该公式说明
在这个维数，同调信息可以给出几何信息。在更高维，类似命题~\ref{prop:9-4-4} 的论证至多
只能确定\BookRecordIndex{T0283}求值映射的度数，而不能证明它是微分同胚。相关结果见 McDuff \srcite{259}。
其中还证明了 Eliashberg--Floer--McDuff 对定理~\ref{thm:9-4-2} 的一个推广：若
$(M,\omega)$ 满足 $\omega$ 在 $\pi_2(M)$ 上消失，并且在某个紧集外辛同胚于欧氏空间，
则 $M$ 可缩。不过，当 $n>2$ 时，仍不知道无穷远附近定义的这个到 $\mathbb R^{2n}$ 的
辛同胚能否延拓成整个 $M$ 上定义的辛同胚。
\end{enumerate}
\end{remark}

说完这些，现在开始证明各定理。先证明定理~\ref{thm:9-4-1}(i)、(ii)，再证明
定理~\ref{thm:9-4-2}，最后证明定理~\ref{thm:9-4-1}(iii)。论证基于两个基本结果。
第一个涉及自交数为 $p$ 的嵌入 $A$-球面\BookRecordIndex{T0246}模空间，其中 $p$ 非负且极小。在这种情形，
命题~\ref{prop:9-4-4} 断言无标记球面的\BookRecordIndex{T0246}模空间 $\mathcal M(A;J)/G$ 紧，并且带适当数目
\BookRecordIndex{T0066}标记点的\BookRecordIndex{T0283}求值映射是微分同胚。由此可知，要么 $p=1$ 且 $M=\mathbb{CP}^2$，要么
$p=0$ 且 $M$ 是直纹曲面。第二个基本结果是定理~\ref{thm:9-4-7}，它细致分析了在集合
\[
  X:=S^2\times\{\mathrm{pt}\}\cup\{\mathrm{pt}\}\times S^2
\]
附近为标准形式的 $S^2\times S^2$ 上\BookRecordIndex{T0363}辛形式。

在下一个命题中，$\Delta^k\subset M^k$ 表示\BookRecordIndex{T0082}大对角线，即所有不两两互异的
$k$-元组 $\mathbf x=(x_1,\ldots,x_k)$ 所成的集合。

\begin{proposition}\label{prop:9-4-4}
设 $(M,\omega)$ 是紧连通辛四维流形，$J$ 是 $M$ 上任意 $\omega$-\BookRecordIndex{T0372}驯服\BookRecordIndex{T0088}殆复结构，而
$A\in H_2(M;\mathbb Z)$ 满足
\[
  A\mathbin{\boldsymbol\cdot}A=:p\ge0.
\]
假设 $A$ 可由辛嵌入二维球面表示，且不存在满足
$-1\le S\mathbin{\boldsymbol\cdot}S<p$ 的辛嵌入二维球面 $S$。再假设每个非常值
$J$-全纯球面的 Chern 数都为正。则表示类 $A$ 的每个 $J$-全纯球面都是嵌入的（因而
正则），\BookRecordIndex{T0246}模空间 $\mathcal M(A;J)/G$ 紧，并且\BookRecordIndex{T0283}求值映射
\begin{equation}\label{eq:9-4-1}
  \operatorname{ev}:\mathcal M_{0,p+1}(A;J)
  \longrightarrow M^{p+1}\setminus\Delta^{p+1}
\end{equation}
是微分同胚。
\end{proposition}

这里的 $J$ 是使每个非常值 $J$-全纯球面都有正 Chern 数的任意\BookRecordIndex{T0372}驯服\BookRecordIndex{T0088}殆复结构；我们不假设
它正则。在四维，这个假设略弱于\BookRecordIndex{T0408}正则性：事实上，引理~\ref{lem:6-4-4} 的证明说明，若
$J\in\mathcal J_{\mathrm{reg}}(M,\omega)$，则每个非常值 $J$-全纯球面
$u:S^2\to M$ 都有正 Chern 数。还要注意，对 $A$ 的假设涉及嵌入 \BookRecordIndex{T0003}$J$-全纯曲线的性质。
因此，需要用\BookRecordIndex{T0055}伴随公式\BookRecordIndex{T0326}推出 $A$ 是 \BookRecordIndex{T0002}$J$-不可分解的；如引理~\ref{lem:7-1-8} 所述，这又
蕴含 $\mathcal M(A;J)/G$ 的\BookRecordIndex{T0192}紧性。我们证明\BookRecordIndex{T0283}求值映射为微分同胚时，要像
引理~\ref{lem:3-3-2} 的\BookRecordIndex{T0408}正则性结果的证明那样使用实线性\BookRecordIndex{T0210}柯西--黎曼算子的
Riemann--Roch 定理。

\begin{corollary}[格罗莫夫]\label{cor:9-4-5}
令 $M:=\mathbb{CP}^2$，并令 $\omega=\omega_{\mathrm{FS}}$ 为标准\BookRecordIndex{T0363}辛形式。则对每个
$J\in\mathcal J_\tau(\mathbb{CP}^2,\omega_{\mathrm{FS}})$ 和任意两个不同的点
$x_1,x_2\in\mathbb{CP}^2$，都存在一个在重参数化意义下唯一的 $J$-全纯球面
$u:S^2\to\mathbb{CP}^2$，满足
$[u]=[\mathbb{CP}^1]$ 且 $x_1,x_2\in u(S^2)$。
\end{corollary}

\begin{proof}
取 $p=1$ 时，命题~\ref{prop:9-4-4} 的假设都成立。特别地，因为
$\mathbb{CP}^2$ 单调，每个非常值 $J$-全纯球面都有正 Chern 数。
\end{proof}

命题~\ref{prop:9-4-4} 的证明以如下观察为基础。

\begin{lemma}\label{lem:9-4-6}
设 $(M,\omega,J)$ 和 $A$ 满足命题~\ref{prop:9-4-4} 的假设，则 $A$ 是
\BookRecordIndex{T0002}$J$-不可分解的。
\end{lemma}

\begin{proof}
反设不然。则存在同调类 $A_i\in H_2(M;\mathbb Z)$ 和正整数 $m_i$，使每个类 $A_i$
都由单纯 $J$-全纯球面 $u_i:S^2\to M$ 表示，并且
\[
  A=\sum_{i=1}^Nm_iA_i,
  \qquad
  \sum_{i=1}^Nm_i>1,
  \qquad
  A_i\ne A_j\quad(i\ne j).
\]
下面证明，对所有 $i,j$，
\begin{equation}\label{eq:9-4-2}
  A_i\mathbin{\boldsymbol\cdot}A_j\ge0,
  \qquad
  A_i\mathbin{\boldsymbol\cdot}A_i\ge c_1(A_i)>0.
\end{equation}
第一个不等式来自\BookRecordIndex{T0186}交点正性（定理~\ref{thm:2-6-3}）。为证明第二个，注意每个 $u_i$
都是非常值 $J$-全纯球面，因而有正 Chern 数；又因 $\sum_i m_i>1$，故对每个 $i$ 都有
\[
  0<c_1(A_i)<c_1(A).
\]
此外，由\BookRecordIndex{T0054}伴随不等式（定理~\ref{thm:2-6-4}），
\[
  A_i\mathbin{\boldsymbol\cdot}A_i\ge c_1(A_i)-2,
\]
且等号成立当且仅当 $u_i$ 是嵌入。若 $u_i$ 不是嵌入，则
$A_i\mathbin{\boldsymbol\cdot}A_i>c_1(A_i)-2$；又因为
$A_i\mathbin{\boldsymbol\cdot}A_i-c_1(A_i)$ 总是偶数，便有
$A_i\mathbin{\boldsymbol\cdot}A_i\ge c_1(A_i)$。所以，对 $u_i$ 非嵌入的每个 $i$，
式~\eqref{eq:9-4-2} 的第二个不等式成立。另一方面，若 $u_i$ 是嵌入，则
$A_i\mathbin{\boldsymbol\cdot}A_i=c_1(A_i)-2$。把它与等式
$A\mathbin{\boldsymbol\cdot}A=c_1(A)-2$ 和不等式
$1\le c_1(A_i)<c_1(A)$ 结合，得到
\[
  -1\le A_i\mathbin{\boldsymbol\cdot}A_i
  <A\mathbin{\boldsymbol\cdot}A,
\]
这被题设排除。至此证明了式~\eqref{eq:9-4-2}。

现在考虑恒等式
\[
\begin{split}
  -2
  &=A\mathbin{\boldsymbol\cdot}A-c_1(A)\\
  &=\sum_{i=1}^N(m_i^2-m_i)
       A_i\mathbin{\boldsymbol\cdot}A_i
    +2\sum_{i<j}m_im_jA_i\mathbin{\boldsymbol\cdot}A_j\\
  &\qquad
    +\sum_{i=1}^Nm_i
       \bigl(A_i\mathbin{\boldsymbol\cdot}A_i-c_1(A_i)\bigr).
\end{split}
\]
由式~\eqref{eq:9-4-2}，右边每一项都非负，矛盾。因此反设不成立，故 $A$ 是
\BookRecordIndex{T0002}$J$-不可分解的。
\end{proof}

\begin{proof}[命题~\ref{prop:9-4-4} 的证明]
先假设\BookRecordIndex{T0246}模空间 $\mathcal M(A;J)$ 非空。因为 $A$ 可由辛嵌入球面表示，
\[
  c_1(A)=A\mathbin{\boldsymbol\cdot}A+2.
\]
故由定理~\ref{thm:2-6-4}，表示类 $A$ 的每个单纯 $J$-全纯球面都是嵌入的。此外，由
引理~\ref{lem:9-4-6}，$A$ 是 \BookRecordIndex{T0002}$J$-不可分解的。因此，类 $A$ 中每个 $J$-全纯球面都
单纯，进而嵌入。由引理~\ref{lem:3-3-3}，对每个 $J$-全纯球面
$u\in\mathcal M(A;J)$，算子 $D_u$ 都满射。于是，由定理~\ref{thm:3-1-6}，\BookRecordIndex{T0246}模空间
$\mathcal M(A;J)$ 是维数
\[
  \dim\mathcal M(A;J)=2p+8
\]
的光滑流形。由于 $A$ 是 \BookRecordIndex{T0002}$J$-不可分解的，定理~\ref{thm:5-3-1} 说明商空间
$\mathcal M(A;J)/G$ 紧。

先证明\BookRecordIndex{T0283}求值映射~\eqref{eq:9-4-1} 单射。若
$\operatorname{ev}(u,\mathbf z)=\operatorname{ev}(u',\mathbf z')$，则曲线 $u,u'$ 在
$p+1$ 个不同点相交。由于 $A\mathbin{\boldsymbol\cdot}A=p$，且两条曲线都是嵌入的，
习题~\ref{ex:2-6-7} 断言 $u(S^2)=u'(S^2)$。因此，由推论~\ref{cor:2-5-4}，存在唯一
\BookRecordIndex{T0040}Möbius 变换
$\varphi:S^2\to S^2$，使 $u\circ\varphi=u'$。因为 $u,u'$ 都是嵌入，便有
$\varphi(z_i')=z_i$，$i=0,\ldots,p$，所以
$[u,\mathbf z]=[u',\mathbf z']$。

证明 $d\widetilde{\operatorname{ev}}$ 处处满射略为困难。为此，注意带点映射\BookRecordIndex{T0246}模空间
$\widetilde{\mathcal M}_{0,p+1}(A;J)$ 在
$(u,\mathbf z)=(u,z_0,\ldots,z_p)$ 处的切空间是
\[
  T_{(u,\mathbf z)}\widetilde{\mathcal M}_{0,p+1}(A;J)
  =\ker D_u\oplus T_{z_0}S^2\oplus\cdots\oplus T_{z_p}S^2,
\]
且\BookRecordIndex{T0283}求值映射的微分各分量为
\[
  d\widetilde{\operatorname{ev}}_i(u,\mathbf z)(\xi,\boldsymbol\zeta)
  =\xi(z_i)+du(z_i)\zeta_i,
\]
其中 $\xi\in\ker D_u$，$\zeta_i\in T_{z_i}S^2$。令
$L_0:=\operatorname{im}du\subset u^*TM$ 为曲线 $\operatorname{im}u$ 的切丛，并用
$L_1\subset u^*TM$ 表示它的正交补。于是 $L_0$ 的 Chern 数为二，$L_1$ 的 Chern 数
为 $p$。此外，算子 $D_u$ 分裂为
\[
  D_u(\xi_0+\xi_1)=D_{00}\xi_0+D_{01}\xi_1+D_{11}\xi_1,
\]
其中
$D_{jj}:\Omega^0(S^2,L_j)\to\Omega^{0,1}(S^2,L_j)$ 是实线性\BookRecordIndex{T0210}柯西--黎曼算子，而
$D_{01}:\Omega^0(S^2,L_1)\to\Omega^{0,1}(S^2,L_0)$ 是零阶算子。

先证明映射
\begin{equation}\label{eq:9-4-3}
  \ker D_{11}\longrightarrow
  L_{1,z_0}\oplus\cdots\oplus L_{1,z_p},
  \qquad
  \xi_1\longmapsto\bigl(\xi_1(z_0),\ldots,\xi_1(z_p)\bigr)
\end{equation}
是双射。首先，式~\eqref{eq:9-4-3} 是两个同为 $2p+2$ 维的向量空间之间的实线性映射，
所以只需证明它单射。否则，存在非零截面 $\xi_1\in\ker D_{11}$，在 $p+1$ 个不同点
消失。然而，由于 $D_{11}$ 是实线性\BookRecordIndex{T0210}柯西--黎曼算子，$\xi_1$ 的每个零点都有正的交指标。
（正如在定理~C.1.10 的证明末尾所见，这是因为 $\xi_1$ 属于某个复线性\BookRecordIndex{T0210}柯西--黎曼算子的
核，因而关于 $L_1$ 上某个合适的复结构全纯。）因为 $L_1\to S^2$ 是次数为 $p$ 的复
线丛，$\xi_1$ 至多有 $p$ 个零点，矛盾。因此式~\eqref{eq:9-4-3} 是双射。

明白这一点后，立刻可知 $d\widetilde{\operatorname{ev}}(u,\mathbf z)$ 的核就是\BookRecordIndex{T0413}重参数化群
$\operatorname{PSL}(2,\mathbb C)$ 在 $\widetilde{\mathcal M}_{0,p+1}(A;J)$ 上作用时，经过
$(u,\mathbf z)$ 的轨道切空间。所以，商空间上 $\operatorname{ev}$ 的微分是单射。由于
$\operatorname{ev}$ 是两个同为 $4p+4$ 维的流形之间的光滑映射，它的微分在每一点都是
双射。

下一步证明\BookRecordIndex{T0283}求值映射~\eqref{eq:9-4-1} 是恰当的。设
$[u^\nu,\mathbf z^\nu]$ 是 $\mathcal M_{0,p+1}(A;J)$ 中的序列，并且
\[
  \mathbf x^\nu:=\operatorname{ev}(u^\nu,\mathbf z^\nu)
  \in M^{p+1}\setminus\Delta^{p+1}
\]
收敛到 $\mathbf x\in M^{p+1}\setminus\Delta^{p+1}$。由于
$\mathcal M(A;J)/G$ 紧，必要时重参数化 $u^\nu$ 后，可假设
$(u^\nu,\mathbf z^\nu)$ 有收敛子列。令
$(u,\mathbf z)\in\mathcal M(A;J)\times(S^2)^{p+1}$ 为其极限。因为
$\mathbf x\in M^{p+1}$ 是一个各点两两不同的元组，
$\mathbf z\in(S^2)^{p+1}$ 也由两两不同的点组成。因此，该子列在
$\mathcal M_{0,p+1}(A;J)$ 中收敛。至此，\BookRecordIndex{T0283}求值映射~\eqref{eq:9-4-1} 已被证明为恰当的
单射，并且其微分在每一点都是双射。这说明，只要 $\mathcal M(A;J)\ne\varnothing$，
式~\eqref{eq:9-4-1} 就是微分同胚。

还需证明，对每个使所有非常值 $J$-全纯球面都有正 Chern 数的 $\omega$-\BookRecordIndex{T0372}驯服 $J$，
$\mathcal M(A;J)$ 都非空。首先注意，$M$ 上的 $\omega$-\BookRecordIndex{T0348}相容\BookRecordIndex{T0088}殆复结构是某个纤维可缩丛
的截面，故可构造 $J\in\mathcal J(M,\omega)$，使 $C$ 为 $J$-全纯的。由
引理~\ref{lem:3-3-3}，$C$ 是正则 \BookRecordIndex{T0003}$J$-全纯曲线，因而在 $J$ 的小扰动下持续存在。
所以，存在邻近的\BookRecordIndex{T0088}殆复结构
$J'\in\mathcal J_{\mathrm{reg}}(M,\omega)$，容许类 $A$ 中的 $J'$-全纯球面。因为
$J'$ 正则，每个非常值 $J'$-全纯球面都有正 Chern 数。上面已经证明，在这些条件下，相应
的\BookRecordIndex{T0283}求值映射~\eqref{eq:9-4-1} 是微分同胚。因此得到非平凡\BookRecordIndex{T0143}格罗莫夫--威滕不变量
\[
  \operatorname{GW}^{M}_{A,p+1}(a,\ldots,a)=\pm1,
  \qquad
  a:=\operatorname{PD}([\mathrm{pt}])\in H^4(M).
\]
这又蕴含：对每个 $J\in\mathcal J_\tau(M,\omega)$，类 $A$ 都可由
$J$-全纯\BookRecordIndex{T0340}稳定映射表示。若每个非常值 $J$-全纯球面都有正 Chern 数，则由
引理~\ref{lem:9-4-6}，这种\BookRecordIndex{T0340}稳定映射只能是嵌入 $J$-全纯球面，因而在这种情形
$\mathcal M(A;J)\ne\varnothing$。命题得证。
\end{proof}

\begin{proof}[定理~\ref{thm:9-4-1}(i)、(ii) 的证明]
设 $C\subset M$ 是自交数
$C\mathbin{\boldsymbol\cdot}C=:p\ge0$ 的辛嵌入二维球面，并假设不存在自交数介于
$-1$ 与 $p-1$ 之间的辛嵌入二维球面。令
$A:=[C]$，并取 $J\in\mathcal J_{\mathrm{reg}}(M,\omega)$。则由
命题~\ref{prop:9-4-4}，每条 $A$-曲线都是嵌入的，且\BookRecordIndex{T0283}求值映射
\[
  \operatorname{ev}:\mathcal M_{0,p+1}(A;J)
  \longrightarrow M^{p+1}\setminus\Delta^{p+1}
\]
是微分同胚。

我们通过证明否则类 $A$ 中会有非嵌入曲线，来证明 $p=0$ 或 $p=1$。设 $p\ge2$。
固定 $p-1$ 个不同的点 $x_0,x_3,\ldots,x_p\in M$，并取
$v_1,v_2\in T_{x_0}M$，使向量 $v_1,Jv_1,v_2,Jv_2$ 构成 $T_{x_0}M$ 的一组基。考虑
序列
\[
  x_1^\nu:=\exp_{x_0}(2^{-\nu}v_1),
  \qquad
  x_2^\nu:=\exp_{x_0}(2^{-\nu}v_2).
\]
因为式~\eqref{eq:9-4-1} 是微分同胚，对每个 $\nu$ 都存在唯一的 $J$-全纯球面
$u^\nu:S^2\to M$ 表示类 $A$，满足
\[
  u^\nu(0)=x_0,
  \qquad
  u^\nu(1)=x_1^\nu,
  \qquad
  u^\nu(2)=x_2^\nu,
  \qquad
  x_i\in u^\nu(S^2)\quad(i\ge3).
\]
由于 $\mathcal M(A;J)/G$ 紧，存在 \BookRecordIndex{T0040}Möbius 变换列 $\varphi^\nu\in G$，使
$\varphi^\nu(0)=0$，且 $u^\nu\circ\varphi^\nu$ 有收敛子列。取此子列后，不妨设
\[
  u:=\lim_{\nu\to\infty}u^\nu\circ\varphi^\nu
\]
在 $C^\infty$ 拓扑下存在。再取子列，可假设极限
\[
  z_1:=\lim_{\nu\to\infty}(\varphi^\nu)^{-1}(1),
  \qquad
  z_2:=\lim_{\nu\to\infty}(\varphi^\nu)^{-1}(2)
\]
存在。于是
\[
  u(z_1)
  =\lim_{\nu\to\infty}u^\nu\circ\varphi^\nu
       \bigl((\varphi^\nu)^{-1}(1)\bigr)
  =\lim_{\nu\to\infty}u^\nu(1)
  =\lim_{\nu\to\infty}x_1^\nu=x_0,
\]
同理 $u(z_2)=x_0$。断言 $z_1,z_2$ 不可能都等于零。事实上，若
$z_1^\nu:=(\varphi^\nu)^{-1}(1)\to0$，再取子列后，可假设商
$z_1^\nu/|z_1^\nu|$ 在 $S^1$ 中收敛。考察局部坐标中的差商
\[
  \frac{u^\nu\circ\varphi^\nu(z_1^\nu)
        -u^\nu\circ\varphi^\nu(0)}{|z_1^\nu|}
  =\frac{x_1^\nu-x_0}{|z_1^\nu|},
\]
可知向量 $v_1\in T_{x_0}M$ 属于 $du(0)$ 的像。同一论证适用于序列
$z_2^\nu:=(\varphi^\nu)^{-1}(2)$。但按构造，向量 $v_1,v_2$ 复线性无关，故不可能都属于
$du(0)$ 的像。因此 $z_1\ne0$ 或 $z_2\ne0$。无论哪种情形，$u$ 都有一个自交点；这与
表示类 $A$ 的所有 $J$-全纯球面均为嵌入矛盾。因此 $p\ge2$ 的假设错误，故 $p\le1$，
这证明了 (i)。

下面证明 (ii)。设 $A$ 是辛嵌入球面 $C$ 的同调类，且
$C\mathbin{\boldsymbol\cdot}C=0$；取
$J\in\mathcal J_{\mathrm{reg}}(M,\omega)$。每个非常值 $J$-全纯球面都有正 Chern 数，
所以由命题~\ref{prop:9-4-4} 和定理~\ref{thm:2-6-3}，\BookRecordIndex{T0246}模空间
$\mathcal M(A;J)$ 有如下性质：
\begin{enumerate}[label=(\alph*),leftmargin=3em]
\item 每个 $u\in\mathcal M(A;J)$ 都是嵌入；此外，若
$u_0,u_1\in\mathcal M(A;J)$ 且
$u_0(S^2)\cap u_1(S^2)\ne\varnothing$，则
$u_0(S^2)=u_1(S^2)$；
\item 对每个 $u\in\mathcal M(A;J)$，$D_u$ 都满射；
\item 商空间 $\mathcal M(A;J)/G$ 紧；
\item \BookRecordIndex{T0283}求值映射
$\operatorname{ev}:\mathcal M_{0,1}(A;J)\to M$ 是微分同胚。
\end{enumerate}
现在，$\mathcal M_{0,1}(A;J)$ 是紧定向二维流形 $\mathcal M(A;J)/G$ 上的一个丛，
纤维微分同胚于 $S^2$。此外，拉回形式 $\operatorname{ev}^*\omega$ 在这些纤维上非退化，
因为它们被映为 $M$ 中的 \BookRecordIndex{T0003}$J$-全纯曲线。这证明了定理~\ref{thm:9-4-1}(ii)。
\end{proof}

定理~\ref{thm:9-4-1}(iii) 的证明建立在下面将说明的
$S^2\times S^2$ 上\BookRecordIndex{T0363}辛形式结构定理之上。

\begin{theorem}\label{thm:9-4-7}
设 $(M,\omega)$ 是不含自交数为 $-1$ 的辛嵌入二维球面的紧连通辛四维流形。设
$A,B\in H_2(M;\mathbb Z)$ 满足
\[
  A\mathbin{\boldsymbol\cdot}A=B\mathbin{\boldsymbol\cdot}B=0,
  \qquad
  A\mathbin{\boldsymbol\cdot}B=1,
\]
并假设 $A,B$ 都可由辛嵌入二维球面表示。取面积形式
$\sigma\in\Omega^2(S^2)$，满足 $\int_{S^2}\sigma=1$。则：
\begin{enumerate}[label=(\roman*),leftmargin=3em]
\item 存在微分同胚 $\psi:S^2\times S^2\to M$，使
\[
  \psi^*\omega=a\pi_1^*\sigma+b\pi_2^*\sigma,
  \qquad
  a:=\int_A\omega,
  \quad
  b:=\int_B\omega.
\]
\item 设 $U\subset S^2$ 是开圆盘，而
\[
  \iota:(S^2\times U)\cup(U\times S^2)\longrightarrow M
\]
是嵌入，满足
\[
  \iota^*\omega=a\pi_1^*\sigma+b\pi_2^*\sigma,
  \qquad
  \iota_*[S^2\times\{w\}]=A,
  \quad
  \iota_*[\{z\}\times S^2]=B
\]
对所有 $z,w\in U$ 成立。若 $D\subset U$ 是任意紧集，则 (i) 中的微分同胚 $\psi$ 可
选成在 $(S^2\times D)\cup(D\times S^2)$ 上与 $\iota$ 一致。
\end{enumerate}
\end{theorem}

\begin{proof}
取 $J\in\mathcal J_{\mathrm{reg}}(M,\omega)$。我们在定理~\ref{thm:9-4-1}(ii) 的证明
中已经看到，\BookRecordIndex{T0246}模空间 $\mathcal M(A;J)$ 的元素都是嵌入球面，其像或者不交，或者重合。
\BookRecordIndex{T0246}模空间 $\mathcal M(B;J)$ 同样如此。由\BookRecordIndex{T0186}交点正性（定理~\ref{thm:2-6-3}），每个
$J$-全纯球面 $u\in\mathcal M(A;J)$ 都与每个
$J$-全纯球面 $v\in\mathcal M(B;J)$ 恰好横截相交一次。因此映射
\[
  \Psi:\mathcal M(A;J)/G\times\mathcal M(B;J)/G\longrightarrow M,
  \qquad
  \Psi([u],[v]):=\operatorname{im}[u]\cap\operatorname{im}[v]
\]
良定义。下面证明它是微分同胚。

由取 $p=0$ 的命题~\ref{prop:9-4-4}，\BookRecordIndex{T0283}求值映射
\[
  \operatorname{ev}_A:\mathcal M_{0,1}(A;J)\to M,
  \qquad
  \operatorname{ev}_B:\mathcal M_{0,1}(B;J)\to M
\]
都是微分同胚。所以，通过 $M$ 中每一点都恰有一条 $A$-曲线和一条 $B$-曲线。因此
$\Psi$ 是双射，而其逆可表示为复合
\[
  \Psi^{-1}
  =(\pi_A\times\pi_B)\circ
    (\operatorname{ev}_A\times\operatorname{ev}_B)^{-1}\circ\iota_\Delta,
\]
其中
\[
  \pi_A:\mathcal M_{0,1}(A;J)\to\mathcal M(A;J)/G,
  \qquad
  \pi_B:\mathcal M_{0,1}(B;J)\to\mathcal M(B;J)/G
\]
是显然的投影，$\iota_\Delta:M\to M\times M$ 是对角线的包含映射。故
$\Psi^{-1}$ 光滑。此外，在满足 $u(z)=v(w)$ 的点
$([u,z],[v,w])$ 处，$\operatorname{ev}_A\times\operatorname{ev}_B$ 的微分把竖直空间
映到子空间 $\operatorname{im}du(z)\times\operatorname{im}dv(w)$，而这个子空间横截于
对角线。因此，$(\operatorname{ev}_A\times\operatorname{ev}_B)^{-1}$ 的微分把对角线
的切空间映到一个水平空间。于是 $\Psi^{-1}$ 的微分在每一点都是双射，故 $\Psi$ 是
微分同胚。

现在固定 $u_*\in\mathcal M(A;J)$、$v_*\in\mathcal M(B;J)$。因为 $\Psi$ 是微分同胚，
它在子集
$\mathcal M(A;J)/G\times\{[v_*]\}$ 上的限制诱导到 $v_*$ 的像上的微分同胚
\[
  g:\mathcal M(A;J)/G\longrightarrow\operatorname{im}v_*;
\]
它把每条无参数曲线 $[u]$ 映为该曲线与 $\operatorname{im}v_*$ 的交点。把 $g$ 的逆与
$v_*$ 预复合，得到微分同胚
\[
  g^{-1}\circ v_*:S^2\longrightarrow\mathcal M(A;J)/G,
  \qquad
  w\longmapsto[u_w],
\]
其中 $[u_w]$ 是通过 $v_*(w)$ 的 $\mathcal M(A;J)/G$ 中唯一元素。类似地，有微分同胚
$S^2\to\mathcal M(B;J)/G$，它把 $z$ 映到唯一的
$[v_z]\in\mathcal M(B;J)/G$，其像在点 $u_*(z)$ 与 $u_*$ 的像相交。将微分同胚
$(z,w)\mapsto([u_w],[v_z])$ 与 $\Psi$ 复合，得到微分同胚
\begin{equation}\label{eq:9-4-4}
  \psi=\psi_J:S^2\times S^2\longrightarrow M,
  \qquad
  \psi(z,w):=\operatorname{im}[u_w]\cap\operatorname{im}[v_z].
\end{equation}
对每个 $w\in S^2$，$S^2\times\{w\}$ 在 $\psi$ 下的像是表示类 $A$ 的嵌入
$J$-全纯球面 $\operatorname{im}[u_w]$。所以，对每个 $w\in S^2$，
$S^2\times\{w\}$ 是 $S^2\times S^2$ 中的 $\psi^*J$-全纯球面，且关于
$\psi^*\omega$ 的面积为 $a$。同样，对每个 $z\in S^2$，子流形
$\{z\}\times S^2$ 是面积为 $b$ 的 $\psi^*J$-全纯曲线。因此，$\psi^*\omega$ 与
形式
\[
  \omega_0:=a\pi_1^*\sigma+b\pi_2^*\sigma
\]
上同调\sourcefootnote{原书把二形式 $\psi^*\omega$ 与 $\omega_0$ 称为
``homologous''；Moser 同痕所需且由前述积分确定的是二者的上同调类相同。}。此外，
$T(S^2\times S^2)$ 的竖直子丛和水平子丛关于 $\omega_0$ 互为辛正交补，并且关于
$\psi^*J$ 都是复子丛。因此 $\omega_0$ \BookRecordIndex{T0372}驯服 $\psi^*J$。因为 $\psi^*\omega$ 也\BookRecordIndex{T0372}驯服
$\psi^*J$，二形式
\[
  \omega_t:=\omega_0+t(\psi^*\omega-\omega_0)
\]
对每个 $t\in[0,1]$ 都是\BookRecordIndex{T0363}辛形式。（也可像稍后那样援引习题~\ref{ex:9-4-8} 来说明这
一步。）Moser 同痕论证（见下面的注~\ref{rem:9-4-9}）说明，存在微分同胚
$f:S^2\times S^2\to S^2\times S^2$，使
\[
  f^*\psi^*\omega=\omega_0.
\]
所以复合 $\psi\circ f:S^2\times S^2\to M$ 是辛同胚。这证明了 (i)。

为证明 (ii)，需要更仔细地选取 $J$。必要时缩小 $U$，可假设辛嵌入 $\iota$ 光滑延拓到
$(S^2\times U)\cup(U\times S^2)$ 的闭包的某个开邻域。用 $K\subset M$ 表示该闭包在
$\iota$ 下的像，并选取\BookRecordIndex{T0088}殆复结构 $J\in\mathcal J(M,\omega)$，使拉回结构
$\iota^*J$ 在子集 $(S^2\times U)\cup(U\times S^2)$ 上是标准乘积结构 $J_0$。断言还可
假设每个非常值 $J$-全纯球面都有正 Chern 数。事实上，像完全含于 $K$ 中的每个
$J$-全纯球面都拉回为 $S^2\times S^2$ 中关于标准复结构的全纯球面，因此只要非常值，
它就有正 Chern 数。这表明，每个 Chern 数非正的非常值全纯球面都必须与 $K$ 的补集相交。
由于维数原因，这些球面不可能正则，故可通过在 $M\setminus K$ 中扰动 $J$ 将其排除；
比较注~\ref{rem:3-2-3}。

现在固定 $z_*,w_*\in D$，并令 $u_*,v_*$ 分别等于 $\iota$ 在球面
$S^2\times\{w_*\}$、$\{z_*\}\times S^2$ 上的限制，即
\[
  u_*(z):=\iota(z,w_*),
  \qquad
  v_*(w):=\iota(z_*,w).
\]
令 $\psi:S^2\times S^2\to M$ 为式~\eqref{eq:9-4-4} 定义的微分同胚。则
$\psi(z,w)$ 是通过 $\iota(z_*,w)$ 的 $A$-球面 $[u_w]$ 与通过
$\iota(z,w_*)$ 的 $B$-球面 $[v_z]$ 的唯一交点。因此 $\psi$ 有如下性质：
\begin{enumerate}[label=(\alph*),leftmargin=3em]
\item 若 $w\in U$，则 $\psi(z,w)\in\iota(S^2\times\{w\})$；若 $z\in U$，则
$\psi(z,w)\in\iota(\{z\}\times S^2)$。此外，对所有 $z,w\in S^2$，
\[
  \psi(z,w_*)=\iota(z,w_*),
  \qquad
  \psi(z_*,w)=\iota(z_*,w).
\]
\item 对所有 $z,w\in S^2$，子流形 $S^2\times\{w\}$ 和 $\{z\}\times S^2$ 关于
$\psi^*\omega$ 都是辛的。
\end{enumerate}
由 (a)，$\psi$ 在集合
\[
  X:=S^2\times\{w_*\}\cup\{z_*\}\times S^2
\]
以及 $U\times U$ 上与 $\iota$ 一致。现在的策略是把 $\psi$ 改造成微分同胚 $\psi'$，
使它在 $(S^2\setminus D)\times D$ 和 $D\times(S^2\setminus D)$ 上与 $\iota$ 一致，
并仍满足 (b)。一旦做到这一点，余下证明便和前面一样。事实上，由下面的
习题~\ref{ex:9-4-8}，
\[
  \omega_0\wedge(\psi')^*\omega>0,
\]
所以线性路径
\[
  \omega_t:=(1-t)\omega_0+t(\psi')^*\omega,
  \qquad0\le t\le1,
\]
仍由\BookRecordIndex{T0363}辛形式组成。由构造，$\omega_t=\omega_0$ 在
$(S^2\setminus D)\times D\cup D\times(S^2\setminus D)$ 上成立，故 Moser 论证给出支集
包含于 $(S^2\setminus D)\times(S^2\setminus D)$ 的微分同胚 $f$，满足
$f^*(\psi')^*\omega=\omega_0$；见注~\ref{rem:9-4-9}。于是 $\psi'\circ f$ 满足全部
要求。

因此只需构造 $\psi'$。令它在两个集合
$(S^2\setminus U)\times U$ 和 $U\times(S^2\setminus U)$ 之外等于 $\psi$。这两个集合
不交，所以\BookRecordIndex{T0213}可分别处理。因而只需在 $(S^2\setminus U)\times U$ 上构造 $\psi'$。
当 $w\in U$ 时，像曲线 $\psi(S^2\times\{w\})$ 与
$\iota(S^2\times\{w\})$ 相等，故可写成
\[
  \psi(z,w)=\iota(\varphi(z,w),w),
  \qquad z\in S^2,\quad w\in U.
\]
当 $z\in U$ 或 $w=w_*$ 时，$\varphi(z,w)=z$。\footnote{若 $d\psi$ 在球面
$w=w_*$ 上也等于恒等映射的微分，那么可以把 $\psi$ 与形如
$(z,w)\mapsto(g_t(z,w),w)$ 的 $C^1$-小同痕复合来定义 $\psi'$。然而实际未必如此；修正
同痕关于 $C^1$-范数可能必须很大。二维辛丛纤维的拉直问题曾在若干场合出现。此处给出的
论证利用了 $J$ 在 $(S^2\setminus U)\times U$ 上可积这一特殊事实，但这并非必要；
McDuff \srcite{262} 引理 3.11 中有更一般的版本。}此外，对每个 $z\in S^2$，球面
$\{z\}\times S^2$ 在 $\psi$ 下的像是 $B$-曲线 $\operatorname{im}[v_z]$，故
\[
  \iota^{-1}(\operatorname{im}[v_z])\cap(S^2\times U)
  =\bigl\{(\varphi(z,w),w)\mid w\in U\bigr\}.
\]
因为 $\iota^*J$ 在 $S^2\times U$ 上是标准复结构，这个集合是
$S^2\times U$ 的复子流形，所以对每个 $z\in S^2$，映射
\[
  U\longrightarrow S^2,
  \qquad w\longmapsto\varphi(z,w)
\]
关于通常复结构全纯。特别地，它的 Jacobian 行列式处处非负，因而把面积形式拉回为非负
形式。

现在选取光滑映射 $\rho:S^2\to S^2$，使
\[
  \rho(U)\subset U,
  \qquad
  \rho(D)=w_*,
  \qquad
  \rho(w)=w\quad(w\notin U),
  \qquad
  \det(d\rho(w))\ge0\quad(w\in S^2).
\]
例如，把 $U$ 与 $\mathbb C$ 中的单位圆盘识别后，可取
$\rho(w)=\beta(|w|)w$，其中 $\beta$ 是单调\BookRecordIndex{T0188}截断函数。于是断言
\[
  \psi'(z,w):=\iota(\varphi(z,\rho(w)),w),
  \qquad z\in S^2,\quad w\in U,
\]
在 $U\times U$ 上与 $\psi$ 一致（因为 $z\in U$ 时
$\varphi(z,w)=z$），并且具有所需性质。唯一不显然之处，是对所有 $z$，$\omega$ 在曲线
$w\mapsto\psi'(z,w)$ 的像上非退化。由于
$\iota^*\omega=\omega_0$ 在 $S^2\times U$ 上成立
\sourcefootnote{原书在此印作 $U\times(S^2\setminus U)$；但当前曲线的第二坐标
$w$ 属于 $U$，其像落在 $S^2\times U$，而定理假设正是在这整个区域上给出标准形式。
原印的因子次序与区域不合。}，只需验证 $\omega_0$ 在曲线
\[
  w\longmapsto(\varphi(z,\rho(w)),w)
\]
的像上非退化。这是成立的，因为由构造，该映射第一分量的 Jacobian 行列式非负。
定理得证。
\end{proof}

\begin{exercise}\label{ex:9-4-8}
设 $(V,\omega_0)$ 是四维辛向量空间，$W$ 是 $V$ 的二维辛子空间。用 $W^\perp$ 表示
$W$ 的辛正交补。设 $\omega$ 是 $V$ 上另一个\BookRecordIndex{T0363}辛形式，它在 $W$ 和 $W^\perp$ 上的限制
都是\BookRecordIndex{T0404}正面积形式。证明 $\omega\wedge\omega_0>0$。
\end{exercise}

\begin{remark}[Moser 同痕]\label{rem:9-4-9}
设 $t\mapsto\omega_t$ 是紧\BookRecordIndex{T0358}辛流形 $M$ 上一族光滑且彼此上同调的\BookRecordIndex{T0363}辛形式。则存在光滑同痕
$t\mapsto f_t$，使 $f_0=\operatorname{id}$，且对每个 $t$ 都有
$f_t^*\omega_t=\omega_0$。先选取一形式族 $\beta_t$，使
$d\beta_t=\partial_t\omega_t$。现在设法把 $f_t$ 构造成一族向量场 $X_t$ 对应的同痕。若
这些向量场选成满足
$\iota(X_t)\omega_t+\beta_t=0$，则
\[
  \frac d{dt}f_t^*\omega_t
  =f_t^*\left(\mathcal L_{X_t}\omega_t+\frac d{dt}\omega_t\right)
  =f_t^*d\bigl(\iota(X_t)\omega_t+\beta_t\bigr)=0,
\]
这正是所需结论。

若 $\omega_t$ 在单连通开子集 $U\subset M$ 上不依赖于 $t$，而 $K\subset U$ 是紧集，
则一形式 $\beta_t$ 可选成在 $K$ 上消失。于是每个 $X_t$ 都在 $K$ 上消失，从而每个
$f_t$ 都在 $K$ 上为恒等。注意，只有选定二形式 $\partial_t\omega_t$ 的原函数
$\beta_t$ 之后，同痕才被确定。可以通过选取 $\beta_t$ 为\BookRecordIndex{T0393}余恰形式来消除这种不确定性：
即关于某个度量，令 $\beta_t$ 属于算子
$d^*:\Omega^2(M)\to\Omega^1(M)$ 的像。这样，构造就完全不涉及选择。更多细节见
\srcite{277} 第 3.2 节。
\end{remark}

与其现在完成定理~\ref{thm:9-4-1} 的证明，不如先证明关于欧氏空间上\BookRecordIndex{T0363}辛形式结构的
定理~\ref{thm:9-4-2}，这样更有效率。下一个引理用来控制我们所构造的辛同胚的支集。

\begin{lemma}\label{lem:9-4-10}
设 $V\subset\mathbb R^{2n}$ 是星形紧集，并假设
$f:\mathbb R^{2n}\setminus V\to\mathbb R^{2n}$ 是在无穷远附近等于恒等映射的辛嵌入。
则对 $V$ 的每个开邻域 $W\subset\mathbb R^{2n}$，都存在辛同胚
$g:\mathbb R^{2n}\to\mathbb R^{2n}$，使
$g|_{\mathbb R^{2n}\setminus W}=f$。
\end{lemma}

\begin{proof}
把 $V$ 换成它与一个以原点为中心的小球之并，不失一般性可假设
$0\in\operatorname{int}(V)$。于是存在常数 $T>1$，使对每个
$x\in\mathbb R^{2n}\setminus V$ 都有 $f(Tx)=Tx$。定义
\[
  f_t:\mathbb R^{2n}\setminus V\longrightarrow\mathbb R^{2n},
  \qquad
  f_t(x):=t^{-1}f(tx),
  \qquad1\le t\le T.
\]
则 $f_1=f$、$f_T=\operatorname{id}$，并且对 $1\le t\le T$ 都有
$f_t^*\omega_0=\omega_0$
\sourcefootnote{原书先只在 $1\le t\le T$ 定义 $f_t$，随后却把这句话的量词印作
``all $t\ge1$''；此处唯一需要且有定义的范围是 $1\le t\le T$。}。记
\[
\begin{split}
  V_t&:=\bigl\{x\in\mathbb R^{2n}\mid
    tx\notin f(\mathbb R^{2n}\setminus tV)\bigr\},\\
  W_t&:=\bigl\{x\in\mathbb R^{2n}\mid
    tx\notin f(\mathbb R^{2n}\setminus tW)\bigr\}.
\end{split}
\]
之所以需要这些颇为曲折的定义，是因为 $f$ 并非处处有定义。事实上，$V_t$ 是
$\mathbb R^{2n}$ 中以
\[
  f_t(\partial V)=t^{-1}f(t\partial V)
\]
为边界的紧区域\sourcefootnote{原书把这里的环境空间印作 $\mathbb R^4$，但引理及整个
证明都在 $\mathbb R^{2n}$ 中进行。}。所以 $V_t$ 紧，$f_t$ 是从
$\mathbb R^{2n}\setminus V$ 到 $\mathbb R^{2n}\setminus V_t$ 的微分同胚，$W_t$ 是
$V_t$ 的开邻域，而且 $f_t$ 把 $\mathbb R^{2n}\setminus W$ 微分同胚地映到
$\mathbb R^{2n}\setminus W_t$。

用
\[
  \frac d{dt}f_t=X_t\circ f_t
\]
定义向量场
$X_t\in\operatorname{Vect}(\mathbb R^{2n}\setminus V_t)$。因为
$f_t^*\omega_0=\omega_0$，一形式
$\alpha_t:=\iota(X_t)\omega_0$ 对每个 $t$ 都闭。此外，$\alpha_t$ 在无穷远附近消失。
由于 $\mathbb R^{2n}\setminus V_t$ 是补集紧的连通开集，这蕴含 $\alpha_t$ 正合（即使
$n=1$ 也如此）。所以存在光滑哈密顿函数族
$F_t\in C^\infty(\mathbb R^{2n}\setminus V_t)$，使
\[
  \iota(X_t)\omega_0=dF_t,
\]
且 $F_t$ 在无穷远附近消失。现在选取光滑\BookRecordIndex{T0188}截断函数族
$\rho_t:\mathbb R^{2n}\to[0,1]$，使 $\rho_t$ 在 $V_t$ 的某个邻域中等于零，并在
$\mathbb R^{2n}\setminus W_t$ 上等于一。于是函数
\[
  G_t:=\rho_tF_t:\mathbb R^{2n}\longrightarrow\mathbb R
\]
在整个 $\mathbb R^{2n}$ 上良定义；它在 $V_t$ 上等于零，并且有（关于 $t$ 一致的）紧
支集。令 $g_t$ 是 $G_t$ 生成的\BookRecordIndex{T0158}哈密顿同痕，即
\[
  \frac d{dt}g_t=Y_t\circ g_t,
  \qquad
  g_T=\operatorname{id},
  \qquad
  \iota(Y_t)\omega_0=dG_t.
\]
则 $g:=g_1$ 是 $\mathbb R^{2n}$ 的辛同胚，并且在
$\mathbb R^{2n}\setminus W$ 上与 $f$ 一致。
\end{proof}

在下面的证明中，记
\[
  D_r:=\{z\in\mathbb C\mid |z|\le r\}
\]
为 $\mathbb C$ 中半径为 $r$ 的闭圆盘，并记
$\Delta_r:=D_r\times D_r$ 为
$\mathbb C^2=\mathbb R^4$ 中半径为 $r$ 的闭\BookRecordIndex{T0110}多圆盘。

\begin{proof}[定理~\ref{thm:9-4-2} 的证明]
取正实数 $R,r,\varepsilon$，使 $R-\varepsilon>r$ 且 $V\subset\Delta_r$。记
\[
  M_R:=\psi(\operatorname{int}(\Delta_R)\setminus V)\cup K
  =M\setminus\psi(\mathbb R^4\setminus\operatorname{int}(\Delta_R)).
\]
（若 $\psi$ 处处有定义，$M_R$ 就只是 $\operatorname{int}(\Delta_R)$ 在 $\psi$ 下的像。）
现在把视作 $S^2\times S^2$ 子集的内部
$\operatorname{int}(\Delta_R)$ 替换为 $M_R$，以构造紧\BookRecordIndex{T0358}辛流形 $\widetilde M$。更明确地，
选取面积形式 $\sigma\in\Omega^2(S^2)$，使
$\int_{S^2}\sigma=\pi R^2$；取保面积嵌入
$\iota:\operatorname{int}(D_R)\to S^2\setminus\{z_*\}$，并定义
\[
  \widetilde M:=
  \bigl((S^2\times S^2)\setminus(\iota\times\iota)(V)\bigr)
  \cup M_R\big/\sim.
\]
这里的等价关系经由
\[
  (\iota(z_1),\iota(z_2))\sim\psi(z_1,z_2),
  \qquad
  (z_1,z_2)\in\operatorname{int}(\Delta_R)\setminus V,
\]
把 $M_R\setminus K$ 与
$(\iota\times\iota)(\operatorname{int}(\Delta_R)\setminus V)$ 识别。因为
$\iota\times\iota$ 和 $\psi$ 都是辛同胚，$\widetilde M$ 上存在唯一\BookRecordIndex{T0363}辛形式
$\widetilde\omega$，它在 $M_R\setminus K$ 上与 $\omega$ 一致，并在
$(S^2\times S^2)\setminus(\iota\times\iota)(V)$ 上与
$\pi_1^*\sigma+\pi_2^*\sigma$ 一致。

由定理~\ref{thm:9-4-7}(ii)，存在微分同胚
$\widetilde\psi:S^2\times S^2\to\widetilde M$，它在
\[
  (S^2\times S^2)\setminus(\iota\times\iota)(\Delta_{R-\varepsilon})
\]
上等于恒等映射，并满足
\[
  \widetilde\psi^*\widetilde\omega
  =\pi_1^*\sigma+\pi_2^*\sigma.
\]
定义映射 $\psi':\mathbb R^4\to M$ 为
\[
  \psi'(z_1,z_2):=
  \begin{cases}
    \widetilde\psi(\iota(z_1),\iota(z_2)),
      &(z_1,z_2)\in\operatorname{int}(\Delta_R),\\
    \psi(z_1,z_2),&(z_1,z_2)\notin\Delta_{R-\varepsilon}.
  \end{cases}
\]
则 $\psi'$ 是良定义的辛同胚，并在
$\mathbb R^4\setminus\Delta_{R-\varepsilon}$ 上与原映射 $\psi$ 一致。

现在令 $U$ 是 $K$ 的开邻域，并记
\[
  W:=V\cup\psi^{-1}(U\setminus K),
\]
于是
\[
  \psi(\mathbb R^4\setminus W)=M\setminus U.
\]
考虑辛嵌入
\[
  f:=(\psi')^{-1}\circ\psi:
  \mathbb R^4\setminus V\longrightarrow\mathbb R^4.
\]
因为 $V$ 星形，由引理~\ref{lem:9-4-10}，存在辛同胚
$g:\mathbb R^4\to\mathbb R^4$，使
$g|_{\mathbb R^4\setminus W}=f$。定义
$\psi'':=\psi'\circ g:\mathbb R^4\to M$。则 $\psi''$ 是辛同胚，它在
$\mathbb R^4\setminus W$ 上与 $\psi$ 一致，并把
$\mathbb R^4\setminus W$ 映到 $M\setminus U$ 上。定理得证。
\end{proof}

有了这个结果，现在可以证明定理~\ref{thm:9-4-1} 的最后一部分。

\begin{proof}[定理~\ref{thm:9-4-1}(iii) 的证明]
设 $C\subset M$ 是微分同胚于二维球面、自交数
$C\mathbin{\boldsymbol\cdot}C=1$，且面积
$\int_C\omega=\pi$ 的辛子流形。辛邻域定理断言，存在
$\mathbb{CP}^1$ 在 $\mathbb{CP}^2$ 中的开邻域 $U$ 以及嵌入
$\iota:U\to M$，使
\[
  \iota^*\omega=\omega_{\mathrm{FS}},
  \qquad
  \iota(\mathbb{CP}^1)=C;
\]
比较 \srcite{277} 的定理 3.30。这里 $\omega_{\mathrm{FS}}$ 是
$\mathbb{CP}^2$ 上使 $\mathbb{CP}^1$ 面积为 $\pi$ 的 Fubini--Study 形式。现在，
$\mathbb{CP}^2\setminus\mathbb{CP}^1$ 辛同胚于 $\mathbb R^4$ 中的开单位球；见下面的
习题~\ref{ex:9-4-11}。任选这样的辛同胚
\[
  \psi:\operatorname{int}(B_1)
  \longrightarrow\mathbb{CP}^2\setminus\mathbb{CP}^1,
  \qquad
  \psi^*\omega_{\mathrm{FS}}=\omega_0.
\]
选取常数 $\varepsilon>0$，使
$\psi(\operatorname{int}(B_1)\setminus B_{1-2\varepsilon})\subset U$，并考虑\BookRecordIndex{T0358}辛流形
\[
  \widetilde M:=(M\setminus C)\cup(\mathbb R^4\setminus B_{1-2\varepsilon})\big/\sim,
\]
其中等价关系为
\[
  x\sim z
  \quad\Longleftrightarrow\quad
  x=\iota(\psi(z)),
  \qquad
  x\in M\setminus C,
  \quad
  z\in\operatorname{int}(B_1)\setminus B_{1-2\varepsilon}.
\]
由定理~\ref{thm:9-4-2}，存在辛同胚
$f:\mathbb R^4\to\widetilde M$，它在
$\mathbb R^4\setminus B_{1-\varepsilon}$ 上等于恒等映射。因此，$f$ 限制为从
$\operatorname{int}(B_1)$ 到 $M\setminus C$ 的辛同胚，并在
$\operatorname{int}(B_1)\setminus B_{1-\varepsilon}$ 上等于 $\iota\circ\psi$。于是辛同胚
\[
  M\longrightarrow\mathbb{CP}^2,
  \qquad
  x\longmapsto
  \begin{cases}
    \psi(f^{-1}(x)),&x\in M\setminus C,\\
    \iota^{-1}(x),&x\in M\setminus f(B_{1-\varepsilon})
  \end{cases}
\]
满足定理要求。
\end{proof}

\begin{exercise}\label{ex:9-4-11}
通过嵌入
\[
  \mathbb{CP}^1\longrightarrow\mathbb{CP}^2,
  \qquad
  [z_0:z_1]\longmapsto[0:z_0:z_1]
\]
的像来识别 $\mathbb{CP}^1$。证明：辛同胚
$\psi:\operatorname{int}(B_1)\to\mathbb{CP}^2\setminus\mathbb{CP}^1$ 的一个显式例子，
是微分同胚
\[
\begin{split}
  \operatorname{int}(B_1)&\longrightarrow\mathbb C^2,\\
  (\zeta_1,\zeta_2)&\longmapsto
  \left(
    \frac{\zeta_1}{\sqrt{1-|\zeta_1|^2-|\zeta_2|^2}},
    \frac{\zeta_2}{\sqrt{1-|\zeta_1|^2-|\zeta_2|^2}}
  \right)
\end{split}
\]
与嵌入
\[
  \mathbb C^2\longrightarrow\mathbb{CP}^2,
  \qquad
  (z_1,z_2)\longmapsto[1:z_1:z_2]
\]
的复合。

\emph{提示：}使用命题~\ref{prop:9-3-3} 的证明中给出的
$\omega_{\mathrm{FS}}$ 公式。
\end{exercise}

\begin{remark}\label{rem:9-4-12}
在高于四维的情形，仅仅存在 $J$-全纯球面并不能给出信息；因为由格罗莫夫的辛嵌入定理，
这种球面既太常见，又会在扰动 $J$ 时消失。见 Gromov \srcite{161} 的
第 3.4(B) 节和 Eliashberg--Mishachev \srcite{102}。一个非平凡\BookRecordIndex{T0143}格罗莫夫--威滕不变量的
存在性更有意义，不过迄今人们对这个条件的几何和动力学后果还远未完全理解。目前大多数
应用都来自小\BookRecordIndex{T0229}量子上同调环的性质；正如第~11.1 节将看到的，这个环由亏格零三点
\BookRecordIndex{T0143}格罗莫夫--威滕不变量得到。

另一个可能有趣的条件，是如下形式的\BookRecordIndex{T0143}格罗莫夫--威滕不变量非零：
\[
  \operatorname{GW}^{M}_{A,k}
  \bigl(\operatorname{PD}(\mathrm{pt}),a_2,\ldots,a_k\bigr)\ne0.
\]
具有这种性质的流形称为\emph{\BookRecordIndex{T0354}辛单有理的}。这一条件有多种几何后果。例如，像第~\ref{sec:9-3} 节
那样可知，一条单有理化曲线所代表之类 $A$ 的辛面积 $\omega(A)$ 会限制辛嵌入球的大小；又如
Hofer--Viterbo \srcite{182} 和 Lu \srcite{253} 所示，单有理条件有助于解决
注~\ref{rem:9-1-3}(iii) 中讨论的 Weinstein 猜想。Hu--Li--Ruan \srcite{192} 证明，
单有理性质在
吹起和吹下下不变，所以用他们的术语说，它是一个双有理辛不变量。

一个有趣的公开问题是：每个单调\BookRecordIndex{T0358}辛流形是否都单有理，就像 \BookRecordIndex{T0207}凯勒 流形那样？（这里，若
不假设 $[\omega]=c_1$ 在 $H_2(M)$ 而非只在 $\pi_2(M)$ 上成立，还应假设 $M$ 单连通。）
另一个问题是理解何种几何条件会迫使这样的\BookRecordIndex{T0143}格罗莫夫--威滕不变量非零。例如，对每个
$\omega$-\BookRecordIndex{T0372}驯服的 $J$，若通过每个点都有一个 $J$-全纯球面，这是否足够？由
定理~\ref{thm:9-4-1} 及其对吹起的推广，这在四维成立。在高维，唯一相关的结果涉及带
$S^1$-作用的流形：McDuff 在 \srcite{273} 中证明它们单有理。Kollár 和 Ruan
\srcite{214,192} 证明，每个单有理光滑射影簇都\BookRecordIndex{T0354}辛单有理。然而，对\emph{\BookRecordIndex{T0385}有理连通}性质
尚不知道类似结果；在辛语境中，这意味着某个不变量
\[
  \operatorname{GW}^{M}_{A,k}
  \bigl(\operatorname{PD}(\mathrm{pt}),\operatorname{PD}(\mathrm{pt}),
  a_3,\ldots,a_k\bigr)
\]
非零。关于这类问题近期一些非常有趣的工作，见 Voisin \srcite{410} 和
Z.~Tian \srcite{396}。
\end{remark}
